\documentclass{beamer}
\usefonttheme{professionalfonts}
%\usetheme{Rochester}
%\usetheme{AnnArbor}
\usetheme{Szeged}
\useoutertheme{infolines}
\usecolortheme{dolphin}
\usepackage{graphicx}
% \usepackage[english]{babel}
%\usepackage{hyperref}
\usepackage{multicol}
\usepackage{color}
\usepackage{xcolor}
%\usepackage{smartdiagram}

\usepackage[utf8]{inputenc}
  \usepackage{listings}




%\definecolor{carmine}{rgb}{0.59, 0.0, 0.09}

\usepackage{graphicx}
\usepackage{tikz}
\usetikzlibrary{shapes.callouts}
\usepackage{multirow}

\usepackage{enumerate}
\usepackage[shortlabels]{enumitem}
\usepackage{amsmath,amsthm,amssymb,amsfonts,mathrsfs,latexsym,stmaryrd}

%\usepackage{amssymb, amsmath, amsbsy}
\usepackage{adjustbox,lipsum}
\usepackage{placeins}

\usepackage{textpos}
\usepackage{wrapfig}
\usepackage[spanish, es-tabla]{babel}

\usepackage{caption}
\setlength{\parindent}{0pt}
%\usepackage[hyperfootnotes=false]{hyperref}

\usepackage{url}
\usepackage{float}
\usepackage{tikz}
\usetikzlibrary{patterns}

\usepackage{lscape}
\usepackage{paralist} 
\usepackage{indentfirst}
\usepackage[dvipsnames,svgnames]{xcolor} 
\usepackage{setspace}
\usepackage{multicol} 
\usepackage{float}

\definecolor{codegreen}{rgb}{0,0.6,0}
\definecolor{codegray}{rgb}{0.5,0.5,0.5}
\definecolor{codepurple}{cmyk}{0.65, 1, 0, 0.35}
\definecolor{backcolour}{rgb}{0.95,0.95,0.92}

\usepackage{listings}

\lstset{
    language=R,
    basicstyle=\small\ttfamily,
    commentstyle=\color{gray},
    keywordstyle=\color{blue},
    stringstyle=\color{black},
    showstringspaces=false,
    tabsize=4,
    breaklines=true,
    postbreak=\mbox{\textcolor{red}{$\hookrightarrow$}\space},
    numbers=left,
    numberstyle=\tiny\color{gray},
    frame=single,
    backgroundcolor=\color{gray!10}
}


\title[ICS-162]{Econometría ICS-294\\ \small{Clase 7: Propiedades estadísticas del estimador MCO}}

\author{Jocelyn Tapia Stefanoni}
\institute[USM]
{Universidad Técnica Federico Santa María\\ % Your institution for the title page
	\medskip
	\textit{} % Your email address
	
}
\logo{

\includegraphics[width=3cm]{figuras/logoUSM.png}
}
\date{}
\newcommand\Fontvi{\fontsize{9}{7.2}\selectfont}
\begin{document}
	
\begin{frame}
	\titlepage % Print the title page as the first slide


\end{frame}
%%%%%%%%%%%%%%%%%%%%%%%%%%%%%%%%%%%%%%%%%%%%%%%%%%%%%%%%%%%%%%%%%%%%%%%%%%%
\logo{}

\begin{frame}[shrink=10]

\frametitle{Resumen clase anterior}

\begin{itemize}
    \item Vimos la forma matricial del modelo de regresión lineal múltiple: $y = X\beta + u$
    \item Demostramos que por MCO se puede obtener: $\hat{\beta} = \begin{bmatrix} \hat{\beta}_0 \\ \hat{\beta}_1 \\ \hat{\beta}_2 \\ \vdots \\ \hat{\beta}_k \end{bmatrix} = (X^\top X)^{-1}X^\top y$
\end{itemize}

Interpretación $\hat{\beta}_1$: Cuando cambia $x_1$ en una unidad, en promedio $y$ cambia $\hat{\beta}_1$ unidades, manteniendo todas las demás $x$ constantes.

\end{frame}
%%%%%%%%%%%%%%%%%%%%%%%%%%%%%%%%%%%%%%%%%%%%%%%%%%%%%%%%%%%%%%%%%%%%%%%%%%%
\begin{frame}[shrink=10]
\frametitle{Ejemplo: obtener $\hat{\beta} = (X^\top X)^{-1}X^\top y$}

Modelo: $y_i = \beta_0 + \beta_1 x_{i1} + \beta_2 x_{i2} + u_i$, $n=5$.

Base de datos: \begin{center}
\resizebox{0.8\textwidth}{!}{%
\begin{tabular}{c|cccc}

Individuos & unos & $x_1$ (experiencia) & $x_2$ (educación) & $y$ (salario) \\
\hline
Juan & 1 & 12 & 93 & 769 \\
Pedro & 1 & 18 & 119 & 808 \\
Sofia & 1 & 14 & 108 & 825 \\
Caro & 1 & 12 & 96 & 650 \\
Alex & 1 & 11 & 74 & 562 \\

\end{tabular}%
}
\end{center}

\end{frame}

\begin{frame}[shrink=10]
\frametitle{Ejemplo: obtener $\hat{\beta} = (X^\top X)^{-1}X^\top y$}

Modelo: $y_i = \beta_0 + \beta_1 x_{i1} + \beta_2 x_{i2} + u_i$, $n=5$.

Base de datos: \begin{center}
\resizebox{0.8\textwidth}{!}{%
\begin{tabular}{c|cccc}

Individuos & unos & $x_1$ (experiencia) & $x_2$ (educación) & $y$ (salario) \\
\hline
Juan & 1 & 12 & 93 & 769 \\
Pedro & 1 & 18 & 119 & 808 \\
Sofia & 1 & 14 & 108 & 825 \\
Caro & 1 & 12 & 96 & 650 \\
Alex & 1 & 11 & 74 & 562 \\

\end{tabular}%
}
\end{center}

\[
X = \begin{bmatrix}
1 & 12 & 93 \\
1 & 18 & 119 \\
1 & 14 & 108 \\
1 & 12 & 96 \\
1 & 11 & 74 \\
\end{bmatrix} \quad 
X^\top = \begin{bmatrix}
1 & 1 & 1 & 1 & 1 \\
12 & 18 & 14 & 12 & 11 \\
93 & 119 & 108 & 96 & 74 \\
\end{bmatrix}
\]
\[
X^\top X = \begin{bmatrix}
5 & 67 & 490 \\
67 & 929 & 6736 \\
490 & 6736 & 49166 \\
\end{bmatrix}
\]
\end{frame}
%%%%%%%%%%%%%%%%%%%%%%%%%%%%%%%%%%%%%%%%%%%%%%%%%%%%%%%%%%%%%%%%%%%%%%%%%%%
\begin{frame}[shrink=10]
\frametitle{Ejemplo: obtener $\hat{\beta} = (X^\top X)^{-1}X^\top y$}

Modelo: $y_i = \beta_0 + \beta_1 x_{i1} + \beta_2 x_{i2} + u_i$, $n=5$.

Base de datos: \begin{center}
\resizebox{0.8\textwidth}{!}{%
\begin{tabular}{c|cccc}

Individuos & unos & $x_1$ (experiencia) & $x_2$ (educación) & $y$ (salario) \\
\hline
Juan & 1 & 12 & 93 & 769 \\
Pedro & 1 & 18 & 119 & 808 \\
Sofia & 1 & 14 & 108 & 825 \\
Caro & 1 & 12 & 96 & 650 \\
Alex & 1 & 11 & 74 & 562 \\

\end{tabular}%
}
\end{center}

\[
X = \begin{bmatrix}
1 & 12 & 93 \\
1 & 18 & 119 \\
1 & 14 & 108 \\
1 & 12 & 96 \\
1 & 11 & 74 \\
\end{bmatrix} \quad 
X^\top = \begin{bmatrix}
1 & 1 & 1 & 1 & 1 \\
12 & 18 & 14 & 12 & 11 \\
93 & 119 & 108 & 96 & 74 \\
\end{bmatrix}
\]
\[
X^\top X = \begin{bmatrix}
5 & 67 & 490 \\
67 & 929 & 6736 \\
490 & 6736 & 49166 \\
\end{bmatrix}
\quad  (X^\top X)^{-1} =
\begin{bmatrix}
8.80 & 0.19 & -0.11 \\
0.19 & 0.17 & -0.02 \\
-0.11 & -0.02 & 0.00 \\
\end{bmatrix}
\]
\end{frame}
%%%%%%%%%%%%%%%%%%%%%%%%%%%%%%%%%%%%%%%%%%%%%%%%%%%%%%%%%%%%%%%%%%%%%%%%%%%
\begin{frame}
\frametitle{Ejemplo: obtener $\hat{\beta} = (X^\top X)^{-1}X^\top y$}


\[
X^\top y =
\begin{bmatrix}
1 & 1 & 1 & 1 & 1 \\
12 & 18 & 14 & 12 & 11 \\
93 & 119 & 108 & 96 & 74 \\
\end{bmatrix}
\times
\begin{bmatrix}
769 \\
808 \\
825 \\
650 \\
562 \\
\end{bmatrix}
= \begin{bmatrix}
3614 \\
49304 \\
360757 \\
\end{bmatrix} 
\]
\end{frame}
%%%%%%%%%%%%%%%%%%%%%%%%%%%%%%%%%%%%%%%%%%%%%%%%%%%%%%%%%%%%%%%%%%%%%%%%%%%
\begin{frame}
\frametitle{Ejemplo: obtener $\hat{\beta} = (X^\top X)^{-1}X^\top y$}


\[
X^\top y =
\begin{bmatrix}
1 & 1 & 1 & 1 & 1 \\
12 & 18 & 14 & 12 & 11 \\
93 & 119 & 108 & 96 & 74 \\
\end{bmatrix}
\times
\begin{bmatrix}
769 \\
808 \\
825 \\
650 \\
562 \\
\end{bmatrix}
= \begin{bmatrix}
3614 \\
49304 \\
360757 \\
\end{bmatrix} = \]
% 
\end{frame}
%%%%%%%%%%%%%%%%%%%%%%%%%%%%%%%%%%%%%%%%%%%%%%%%%%%%%%%%%%%%%%%%%%%%%%%%%%%
\begin{frame}
\frametitle{Ejemplo: obtener $\hat{\beta} = (X^\top X)^{-1}X^\top y$}

Entonces tenemos: $(X^{\top}X)^{-1} = 
\begin{bmatrix}
8.80 & 0.19 & -0.11 \\
0.19 & 0.17 & -0.02 \\
-0.11 & -0.02 & 0.00 \\
\end{bmatrix}$

$X^{\top}y = 
\begin{bmatrix}
3614 \\
49304 \\
360757 \\
\end{bmatrix}$
\end{frame}
%%%%%%%%%%%%%%%%%%%%%%%%%%%%%%%%%%%%%%%%%%%%%%%%%%%%%%%%%%%%%%%%%%%%%%%%%%%
\begin{frame}
\frametitle{Ejemplo: obtener $\hat{\beta} = (X^\top X)^{-1}X^\top y$}

Entonces tenemos: $(X^{\top}X)^{-1} = 
\begin{bmatrix}
8.80 & 0.19 & -0.11 \\
0.19 & 0.17 & -0.02 \\
-0.11 & -0.02 & 0.00 \\
\end{bmatrix}$

$X^{\top}y = 
\begin{bmatrix}
3614 \\
49304 \\
360757 \\
\end{bmatrix}$

\medskip % Salto de línea pequeño

Multiplicando obtenemos los estimadores:
\medskip
$\hat{\beta} = (X^{\top}X)^{-1}X^{\top}y=
\begin{bmatrix}
\hat{\beta}_0 \\
\hat{\beta}_1 \\
\hat{\beta}_2 \\
\end{bmatrix}=
\begin{bmatrix}
140.59 \\
-16.79 \\
8.24 \\
\end{bmatrix}$\\
\medskip
\end{frame}

%%%%%%%%%%%%%%%%%%%%%%%%%%%%%%%%%%%%%%%%%%%%%%%%%%%%%%%%%%%%%%%%%%%%%%%%%%%
%%%%%%%%%%%%%%%%%%%%%%%%%%%%%%%%%%%%%%%%%%%%%%%%%%%%%%%%%%%%%%%%%%%%%%%%%%%
\begin{frame}
\frametitle{Ejemplo: obtener $\hat{\beta} = (X^\top X)^{-1}X^\top y$}

Entonces tenemos: $(X^{\top}X)^{-1} = 
\begin{bmatrix}
8.80 & 0.19 & -0.11 \\
0.19 & 0.17 & -0.02 \\
-0.11 & -0.02 & 0.00 \\
\end{bmatrix}$

$X^{\top}y = 
\begin{bmatrix}
3614 \\
49304 \\
360757 \\
\end{bmatrix}$

\medskip % Salto de línea pequeño

Multiplicando obtenemos los estimadores:
\medskip
$\hat{\beta} = (X^{\top}X)^{-1}X^{\top}y=
\begin{bmatrix}
\hat{\beta}_0 \\
\hat{\beta}_1 \\
\hat{\beta}_2 \\
\end{bmatrix}=
\begin{bmatrix}
140.59 \\
-16.79 \\
8.24 \\
\end{bmatrix}$\\
\medskip
Si quisieramos la matrix de prediccion de y, simplemente calculamos:
\[\hat{y}=X\hat{\beta}\]
\end{frame}

\begin{frame}{En general}

\[
X^T X =
\begin{pmatrix}
n & \sum_{i=1}^{n} x_{i1} & \sum_{i=1}^{n} x_{i2} & \cdots & \sum_{i=1}^{n} x_{ik} \\
\sum_{i=1}^{n} x_{i1} & \sum_{i=1}^{n} x_{i1}^2 & \sum_{i=1}^{n} x_{i1}x_{i2} & \cdots & \sum_{i=1}^{n} x_{i1}x_{ik} \\
\sum_{i=1}^{n} x_{i2} & \sum_{i=1}^{n} x_{i1}x_{i2} & \sum_{i=1}^{n} x_{i2}^2 & \cdots & \sum_{i=1}^{n} x_{i2}x_{ik} \\
\vdots & \vdots & \vdots & \ddots & \vdots \\
\sum_{i=1}^{n} x_{ik} & \sum_{i=1}^{n} x_{i1}x_{ik} & \sum_{i=1}^{n} x_{i2}x_{ik} & \cdots & \sum_{i=1}^{n} x_{ik}^2
\end{pmatrix}
\]
    \[
 X^Ty=\begin{pmatrix}
 \sum_{i=1}^n y_i \\
 \sum_{i=1}^n x_{i1}y_i \\
 \vdots\\
 \sum_{i=1}^n x_{ik}y_i \\
 \end{pmatrix}
 \]
\end{frame}

\begin{frame}{En R}
    \begin{figure}
        \centering
        \includegraphics[width=0.5\linewidth]{sem5/images/datos_R.png}
       
    \end{figure}

  \end{frame}





\begin{frame}{Regresi\'on }

  
    \begin{figure}
        \centering
        \includegraphics[width=0.7\linewidth]{sem5/images/tabla_reg.png}
    \end{figure}
\end{frame}

%%%%%%%%%%%%%%%%%%%%%%%%%%%%%%%%%%%%%%%%%%%%%%%%%%%%%%%%%%%%%%%%%%%%%%%%%%%
%%%%%%%%%%%%%%%%%%%%%%%%%%%%%%%%%%%%%%%%%%%%%%%%%%%%%%%%%%%%
\begin{frame}
\frametitle{Supuestos Clásicos}
¿$\hat{\beta}$ es un buen estimador para $\beta$? ¿Qué significa que sea bueno?\\
\medskip
{\color{blue}
\textbf{\textit{Supuestos clásicos:}}

\begin{enumerate}
    \item Linealidad en parámetros
    \item Muestra aleatoria
    \item Condición de identificación
    \item Ortogonalidad
    \item Homocedasticidad
\end{enumerate}}
\end{frame}
\begin{frame}
\frametitle{Supuesto 1: Linealidad en Parámetros}
\begin{itemize}
    \item La relación poblacional entre las variables sigue un modelo lineal:
    \[
    y = X \beta + u \quad (1)
    \]
        \[ = \beta_0 + \beta_1 x_1 + \beta_2 x_2 + \ldots + \beta_k x_k + u \quad (2)
    \]
    
    donde $\beta_0, \beta_1, \ldots, \beta_k$ son parámetros desconocidos y $u$ es una variable no observable llamada término de error o perturbación.
    \item Si el modelo analítico no es lineal en los parámetros, hay que linealizarlo, en la medida de lo posible.
    \item Por ejemplo:
    \[
    Y = A K^{\alpha_1} Q^{\alpha_2} L^{\alpha_3}
    \]
    se lineariza como
    \[
    \log Y = \log A + \alpha_1 \log K + \alpha_2 \log Q + \alpha_3 \log L
    \]
\end{itemize}
\end{frame}
\begin{frame}
\frametitle{Supuesto 2: Muestra Aleatoria}
\begin{itemize}
    \item Las $n$ observaciones provienen de una muestra aleatoria de la población,
    \[
    \{(x_{i1}, x_{i2}, \ldots, x_{ik}, y_i) : i = 1, 2, \ldots, n\}
    \]
    \item Esto significa que las observaciones son independientes unas de otras, lo que implica que
    \[
    \text{cov}(u_i, u_j|X) = 0 \quad \forall i \neq j
    \]
    \item En series de tiempo lo anterior implica que los errores no tienen autocorrelación.
\end{itemize}
\end{frame}
\begin{frame}
\frametitle{Supuesto 3: Condición de Identificación (No multicolinealidad perfecta)}
\begin{itemize}
    \item $X$ tiene rango columna completo $k + 1$.
    \item No hay relaciones lineales exactas entre las variables independientes.
    \item Esto implica:
    \begin{itemize}
        \item Las columnas de $X$ son linealmente independientes.
        \item Existe la inversa de $X^{\top}X$.
        \item $n \geq k + 1$.
    \end{itemize}
\end{itemize}
\textbf{Ejemplo de violación a este supuesto}: $x_2 = 2x_1$. La fila 2 es linealmente dependiente de la 1.
\end{frame}
\begin{frame}
\frametitle{Supuesto 4: Ortogonalidad (Exogeneidad)}
\begin{itemize}
    \item $E[u|X] = E(u|x_1, x_2, \ldots, x_k) = 0$.
    \item Esto implica que:
    \begin{itemize}
        \item $E[X^{\top}u] = 0$: la covarianza entre las $x$ y el error es cero.
        \item $E(y|X) = E(X\beta + u|X) = X\beta$.
    \end{itemize}
\end{itemize}
Vamos a ver que este supuesto va a ser el supuesto clave para interpretar causalmente los coeficientes.
\end{frame}

\begin{frame}
\frametitle{Supuesto 4: Ortogonalidad (Exogeneidad)}
\begin{itemize}
    \item $E[u|X] = E(u|x_1, x_2, \ldots , x_k) = 0$.
    \item Esto implica que:
    \begin{itemize}
        \item $E[X^{\top}u] = 0$: la covarianza entre las $x$ y el error es cero.
        \item $E(y|X) = E(X\beta + u|X) = X\beta$.
    \end{itemize}
    \item Vamos a ver que este supuesto va a ser el supuesto clave para interpretar causalmente los coeficientes.
    \item Este supuesto puede no cumplirse por varias razones:
    \begin{itemize}
        \item Omisión de una variable que está correlacionada con alguna $x$.
        \item Errores de medición en los regresores.
    \end{itemize}
    \item Por desgracia, muy pocas veces se sabe con seguridad si el valor promedio de los factores no observables de hecho no está relacionado con las variables explicativas. \textcolor{red}{Pero este es el supuesto crítico.}
\end{itemize}
\end{frame}
\begin{frame}
\frametitle{Supuesto 5: Homocedasticidad y no autocorrelación}
\begin{itemize}
    \item La varianza del error no depende de $X$:
    \[
    \text{Var}(u_i|x_1, \ldots , x_k) = E[u_i^2|x] = \sigma^2
    \]
    \item En caso de que el supuesto no se cumpla, se dice que el modelo presenta heterocedasticidad.
    \item Ejemplo donde no se cumpliría esta condición: si tenemos salario en función de años de escolaridad, la varianza de $u$ es probablemente menor para escolaridades pequeñas.

    \item Con esto, la varianza condicional del vector $u$ queda: $\text{Var}(u|X) = \sigma^2I$.
     \item La covarianza de los errores es cero:
    \[
    \text{Covar}(u_i,u_j|x_1, \ldots , x_k) = 0
    \]
\end{itemize}
\end{frame}
\begin{frame}
\frametitle{Supuesto 5: Homocedasticidad}
\begin{itemize}
    \item Homocedasticidad en un modelo simple (mismo $\sigma$).
    
\begin{figure}
    \centering
    \includegraphics[width=0.8\linewidth]{sem5/images/het.png}
\end{figure}
    
    \item Heterocedasticidad en un modelo simple.

\begin{figure}
    \centering
    \includegraphics[width=0.8\linewidth]{sem5/images/het2.png}
\end{figure}

\end{itemize}
\end{frame}
\begin{frame}
\frametitle{Resumen Supuestos}
\begin{enumerate}
    \item Linealidad en parámetros ($y = X\beta + u$).
    \item Muestra aleatoria.
    \item $X$ tiene rango completo ($k+1$), es decir, todas las variables son linealmente independientes.
    \item Ortogonalidad: $E[u|X] = 0$.
    \item Homocedasticidad y no-autocorrelación: $V(u|X) = \sigma^2I$ 
\end{enumerate}
\end{frame}
\begin{frame}
\frametitle{Propiedades Estadísticas de los Estimadores $\hat{\beta}$}
Bajo los supuestos clásicos, los estimadores MCO cumplen lo siguiente:
\begin{enumerate}
    \item $\hat{\beta}$ es insesgado: $E(\hat{\beta}) = \beta$.
    \item $V(\hat{\beta}|X) = (X^{\top}X)^{-1}\sigma^2$.
    \item Teorema de Gauss-Markov: $\hat{\beta}$ es el estimador lineal insesgado de menor varianza (MELI-BLUE).
\end{enumerate}
\end{frame}
%\begin{frame}
% \frametitle{El Estimador MCO es Insesgado}
% \begin{itemize}
%     \item Teorema: Bajo los supuestos 1 a 4, el estimador MCO de $\beta$ es insesgado.
%     \item Demostración:
%     \begin{align*}
%         \hat{\beta} &= (X^{\top}X)^{-1}X^{\top}y \quad \text{(3)} \\
%         &= (X^{\top}X)^{-1}X^{\top}(X\beta + u) \quad \text{(4)} \\
%         &= \beta + (X^{\top}X)^{-1}X^{\top}u \quad \text{(5)}
%     \end{align*}
% \end{itemize}
% \end{frame}
\begin{frame}
\frametitle{El Estimador MCO es Insesgado}
\begin{itemize}
    \item Teorema: Bajo los supuestos 1 a 4, el estimador MCO de $\beta$ es insesgado.
    \item Demostración:
    \begin{align*}
        \hat{\beta} &= (X^{\top}X)^{-1}X^{\top}y \quad \text{(3)} \\
        &= (X^{\top}X)^{-1}X^{\top}(X\beta + u) \quad \text{(4)} \\
        &= \beta + (X^{\top}X)^{-1}X^{\top}u \quad \text{(5)} \\
        E(\hat{\beta}|X) &= \beta + (X^{\top}X)^{-1}X^{\top}E(u|X) \quad \text{(6)} \\
        &= \beta \quad \text{(7)}
    \end{align*}
    \item Esto implica que el promedio de los estimadores $\hat{\beta}$ obtenidos con todas las muestras aleatorias posibles es igual a $\beta$.
\end{itemize}
\end{frame}
\begin{frame}
\frametitle{Varianza del Estimador MCO}
\begin{itemize}
    \item Teorema: Bajo los supuestos 1 a 5, la varianza del estimador MCO de $\beta$ es $\text{Var}(\hat{\beta}|X) = \sigma^2(X^{\top}X)^{-1}$.
    \item Demostración:
    \begin{align*}
        \text{Var}(\hat{\beta}|X) &= E[(\hat{\beta} - E(\hat{\beta}))(\hat{\beta} - E(\hat{\beta}))^{\top}|X] \quad \text{(8)} \\
        &= E[(\hat{\beta} - \beta)(\hat{\beta} - \beta)^{\top}|X] \quad \text{(9)} \\
        &= E[(X^{\top}X)^{-1}X^{\top}u \, u^{\top}X(X^{\top}X)^{-1}|X] \quad \text{(10)} \\
        &= (X^{\top}X)^{-1}X^{\top}E[u \, u^{\top}|X]X(X^{\top}X)^{-1} \quad \text{(11)} \\
        &= (X^{\top}X)^{-1}X^{\top}\sigma^2IX(X^{\top}X)^{-1} \quad \text{(12)} \\
        &= \sigma^2(X^{\top}X)^{-1}(X^{\top}X)(X^{\top}X)^{-1} \quad \text{(13)} \\
        &= \sigma^2(X^{\top}X)^{-1} \quad \text{(14)}
    \end{align*}
    \item nota:  $\hat{\beta}=\beta+(X^{T}X)^{-1}X^{T}u$
\end{itemize}
\end{frame}






\end{document}
